%%%%%%%%%%%%%%%%%%%%%%%%%%%%%%%%%%%%%%%%%%%%%%%%%%%%%%%%%%%%%%%%%%%%%%%%%%%%%%%%
%2345678901234567890123456789012345678901234567890123456789012345678901234567890
%        1         2         3         4         5         6         7         8
% THESIS CONCLUSIONS
\def\baselinestretch{1}
\chapter{Conclusions and future work}
\label{chap:conclusions}

\def\baselinestretch{1.66}

%In this thesis I developed a unified control architecture for a high-paylod, high-Dof robot, with the aim of exploring the robot capabilities to navigate an environment and avoid obstacle using the redundancy in the kinematic chain.
%The control system used data coming from simulated proximity sensor to estimate the relative position of the robot with respect to different types of obstacle.
%I tested the algorithm in different environment, and implemented two major types of obstacle avoidance algorithm: a \textit{single link} and a \textit{multi link}.
%The experimental result shows how the kinematic chain tends to separate the reaching objective trough the link, were the links closer to the base hadles most of the \textit{positon} error while the ones further down the chain are responsible for the \textit{orientation} of the end effector shown here \ref{fig:robot_effects_graph}.\\
%\noindent The second problem was the discontinuities due to high frequency changes in the position of the \textit{minimum distance point} on the robot body. As the robot moves the minimum distance point changes on the robot point cloud, and this can create high frequency changes in the controlled point to avoid the obstacle. To reduce the effect of this problem for the \textit{single link} obstacle avoidance a hysteresis was introduced, but this was not enough in most cases, as shown in \ref{fig:singlelink_snake_result}.\\ 
%\noindent The hysteresis was able to attenuate the link discontinuities but not the one happening on the same link. On the contrary the \textit{multi link} approach revealed to have better performance in attenuating the \textit{activation discontinuities} as shown in \ref{fig:multilink_snake_result_ExtraTg} and \ref{fig:t1_multi_failure_photo_6lnk}.\\
%\noindent The last point of interest that arose from the test, was the handling of the redundancy, as shown in figure \ref{fig:t1_multi_failure_photo_6lnk} the robot can find itself in complex configuration due to the high number of degree of freedom present in the chain. Possible future work should be focused in reducing the discontinuities in the selection of each \textit{minimum distance point}, possibly doing \textit{spatial averaging} of the distance vector direction for $n$ of the neighbors of the \textit{minimum distance point}. This combined with doing \textit{obstacle recognition} on the point cloud be used to allow the robot to incorporate the orientation information in the obstacle avoidance control.\\

\noindent This thesis illustrates the development of a unified control architecture for a high-payload, high-Degree of Freedom robot manipulator. The primary objective was to explore the robot's ability to navigate cluttered environments, avoiding obstacles making use of the redundancy of the kinematic chain. The control system use data generated from simulated proximity, \textit{Time of Flight} sensors to estimate the relative of each link of the robot from the obstacle in the environment. The control algorithm is based on the \textit{task priority approach} that makes use of the null-space projection to allocate degrees of freedom to the tasks. The \textit{obstacle avoidance} task is placed at the maximum priority since it is a safety objective. The strategy used for the \textit{obstacle avoidance} are based on the \textit{minimum distance} approach, referenced at section 2.4 of Chapter:\ref{chap:second}. Two version were implemented: a \textit{single link} and \textit{multi link} version.\\
\noindent Experimental results demonstrate how the kinematic chain partitions the reaching objective across its links. In particular the links closer to the base handles most of the positional error, whereas links further down the chain manage the orientation of the end-effector(see Figure \ref{fig:robot_effects_graph}).\\
A significant challenge identified was the presence of discontinuities in the \textit{activation function} for the avoidance task, caused by high-frequency changes in the location of the minimum distance point on the robot body. As the robot moves, the minimum distance point shifts across the point cloud which can trigger rapid oscillation in the control signal, also because even small changes in the point position can determine big swing in the desired velocity direction. To mitigate this problem in the single-link approach, a hysteresis loop was introduced. However this proved insufficient, as illustrated in \ref{fig:singlelink_snake_result}. The hysteresis loop was effective in reducing changes across whole links of the robot, but not within the links point. Conversely, the multi-link approach demonstrated superior performance in smoothing these activation discontinuities, as shown in \ref{fig:multilink_snake_result_ExtraTg} and \ref{fig:t1_multi_result_9lnk_LowDelta}. The multi link approach distributes the control effort across multiple links, reducing the sensibility to sudden activation and deactivation of the avoidance task. Finally, the handling of remains a point of interest; as seen in Figure \ref{fig:t1_multi_failure_photo_6lnk},the high number of DoF can lead to complex, sub optimal configuration that can block the robot in place.\\




